\thispagestyle{empty}

{\Huge \textbf{Abstract}}

\vspace{0.6cm}

This dissertation is a self-contained account of volatility modelling, from the classical
foundations of stochastic calculus to the rough models that have displaced them, written so that
the mathematics and the finance are visible to each other at every stage.

Part I develops the machinery and the theory built on it. Chapter \ref{ch:foundations} constructs
the It\^o integral from the ground up, with particular attention to the regularity of Brownian
paths: it is the coexistence of a finite, non-zero quadratic variation with an infinite total
variation that both forces It\^o's construction and, pushed a little further, marks the point at
which it fails. Chapter \ref{ch:finance} builds the financial theory on that foundation ---
contingent claims and arbitrage, the fundamental theorems of asset pricing, the Black--Scholes
framework, and the local and stochastic volatility models devised to repair it --- and identifies
at each step which piece of probability is doing the work.

Part II turns to what happens beyond the semimartingale boundary. Chapter \ref{ch:roughpaths}
presents the theory of rough paths: the empirical evidence that volatility is rough, the precise
sense in which It\^o's theory becomes unavailable when the driving signal has critical exponent
greater than two, and Lyons' resolution of that failure through the extension of a path by its
iterated integrals, together with the signature transform and the rough volatility models it
supports. Chapter \ref{ch:applications} brings the theory to data, where the models are calibrated,
implemented and compared.

The organising thread is a single number. The critical exponent of Brownian motion is exactly two:
Young's integral requires less, It\^o's construction reaches it by spending probabilistic structure,
and rough path theory is what remains when even that is out of reach. The empirical exponent of
volatility lies on the far side of it, and the mathematics of Part II is the price of crossing.

\vspace{0.4cm}

\noindent\textbf{Keywords:} stochastic calculus; It\^o integral; Brownian motion; quadratic
variation; volatility; Black--Scholes; local and stochastic volatility; fractional Brownian motion;
rough paths; rough volatility; calibration.

\vspace{1.2cm}

{\Huge \textbf{Resumen}}

\vspace{0.6cm}

Esta memoria presenta un estudio autocontenido de la modelizaci\'on de la volatilidad, desde los
fundamentos cl\'asicos del c\'alculo estoc\'astico hasta los modelos rugosos que los han desplazado,
redactado de modo que la matem\'atica y las finanzas permanezcan mutuamente visibles en todo momento.

La Parte I desarrolla la maquinaria y la teor\'ia construida sobre ella. El Cap\'itulo
\ref{ch:foundations} construye la integral de It\^o desde el principio, prestando especial atenci\'on
a la regularidad de las trayectorias brownianas: es la coexistencia de una variaci\'on cuadr\'atica
finita y no nula con una variaci\'on total infinita lo que obliga a la construcci\'on de It\^o y lo
que, llevado un poco m\'as lejos, marca el punto en el que dicha construcci\'on deja de ser
aplicable. El Cap\'itulo \ref{ch:finance} levanta sobre ese cimiento la teor\'ia financiera
--- derivados y ausencia de arbitraje, los teoremas fundamentales de valoraci\'on de activos, el
marco de Black--Scholes y los modelos de volatilidad local y estoc\'astica concebidos para
corregirlo --- identificando en cada paso qu\'e resultado probabil\'istico est\'a actuando.

La Parte II se ocupa de lo que ocurre m\'as all\'a de la frontera de las semimartingalas. El
Cap\'itulo \ref{ch:roughpaths} expone la teor\'ia de \emph{rough paths}: la evidencia emp\'irica de
que la volatilidad es rugosa, el sentido preciso en que la teor\'ia de It\^o deja de estar disponible
cuando la se\~nal directora tiene exponente cr\'itico mayor que dos, y la soluci\'on de Lyons
mediante la extensi\'on de una trayectoria por sus integrales iteradas, junto con la transformada
signatura y los modelos de volatilidad rugosa que sustenta. El Cap\'itulo \ref{ch:applications}
lleva la teor\'ia a los datos, calibrando, implementando y comparando los modelos.

El hilo conductor es un \'unico n\'umero. El exponente cr\'itico del movimiento browniano es
exactamente dos: la integral de Young exige menos, la construcci\'on de It\^o lo alcanza a costa de
estructura probabil\'istica, y la teor\'ia de \emph{rough paths} es lo que queda cuando ni siquiera
eso est\'a al alcance. El exponente emp\'irico de la volatilidad se sit\'ua al otro lado, y la
matem\'atica de la Parte II es el precio de cruzarlo.

\vspace{0.4cm}

\noindent\textbf{Palabras clave:} c\'alculo estoc\'astico; integral de It\^o; movimiento browniano;
variaci\'on cuadr\'atica; volatilidad; Black--Scholes; volatilidad local y estoc\'astica; movimiento
browniano fraccionario; \emph{rough paths}; volatilidad rugosa; calibraci\'on.

\newpage
